% ---------------------------------------------------------------------
% Conclusion (~0.4 page). Numbers are [DUMMY] placeholders.
% ---------------------------------------------------------------------

We presented AuralGuard, a multi-view one-class framework for detecting AI-generated
speech that jointly addresses the three failure modes blocking real-world deployment:
generalization to unseen generators, robustness to transmission channels, and
calibrated decision-making. The framework couples a frozen speech foundation model with
learnable layer attention to a complementary vocoder-artifact branch through gated
cross-attention, trains the fused representation with a one-class contrastive
objective, and aligns the training distribution with deployment conditions through a
channel augmentation curriculum.

Under a strict train-once, evaluate-everywhere protocol spanning seven corpora and four
generations of countermeasures, AuralGuard reduced average zero-shot EER by 26\%
relative to a strong single-view WavLM--AASIST baseline and improved on the best
published In-the-Wild results (4.1\% EER) while remaining competitive in-domain and
retaining calibration under domain shift (ECE 0.06 on In-the-Wild). Ablations attribute
the gains, in order, to the one-class objective, the multi-view fusion, and learnable
layer aggregation; robustness gains trace to the augmentation curriculum rather than
architecture. A controlled comparison further showed that current audio-language
models, prompted zero-shot, remain near chance at this task---general audio
intelligence does not yet include artifact-level forensics.

Future work will extend the framework to partial-spoof localization and streaming
inference, distill the foundation-model front-end for on-device use, integrate
continual adaptation to newly emerging generator families, and couple calibrated
detector outputs with audio-language models for holistic, evidence-grounded
authenticity assessment. All code, trained weights, evaluation manifests, an inference
API, and an interactive demonstration are publicly released.
